% Abstract — Version B numbers, frozen at commit 9bdb329e (manifest.frozen.json).
\begin{abstract}
\noindent
Large language models produce fluent output that is occasionally fabricated, and detecting these hallucinations is expensive, opaque, or requires access to model internals in most existing approaches. This paper presents HaluRISC, a lightweight black-box framework that predicts hallucination risk from the question, a reference context, and the candidate answer. The system builds 26 engineered evidence-consistency features across seven groups, covering lexical overlap, named entities, NLI entailment, numeric consistency, hedging, semantic similarity, and length, then feeds them to a tuned XGBoost classifier with calibrated probabilities and SHAP explanations. On the HaluEval QA benchmark (20{,}000 rows, 70/15/15 stratified split), the calibrated model reaches F1 = 0.9846, AUROC = 0.9979, and MCC = 0.9697, with an expected calibration error of 0.0045 before any recalibration. A small-sample recalibration on natural RAGTruth responses cuts the ECE on that corpus from 0.8185 to 0.1335, which shows the value of re-calibrating for the target domain. Zero-shot transfer to natural RAGTruth QA still fails (F1 = 0.3022), an honest limitation the paper discusses in detail. The full analysis runs in about 62~ms per sample on a CPU and costs roughly 100 times less than an LLM-as-judge baseline while reaching a higher F1 on a measured 200-sample set. The complete system is deployed as a FastAPI backend with a Next.js and assistant-ui web dashboard for live demonstration.
\end{abstract}
